\documentclass[a4paper]{article}
\usepackage[spanish]{babel}
\usepackage[utf8]{inputenc}
\usepackage{charter}   % tipografia
\usepackage{graphicx}
%\usepackage{makeidx}

%\usepackage{float}
%\usepackage{amsmath, amsthm, amssymb}
%\usepackage{amsfonts}
%\usepackage{sectsty}
%\usepackage{charter}
%\usepackage{wrapfig}
%\usepackage{listings}
%\lstset{language=C}


\usepackage{color} % para snipets de codigo coloreados
\usepackage{fancybox}  % para el sbox de los snipets de codigo

\definecolor{litegrey}{gray}{0.94}

% \newenvironment{sidebar}{%
% 	\begin{Sbox}\begin{minipage}{.85\textwidth}}%
% 	{\end{minipage}\end{Sbox}%
% 		\begin{center}\setlength{\fboxsep}{6pt}%
% 		\shadowbox{\TheSbox}\end{center}}
% \newenvironment{warning}{%
% 	\begin{Sbox}\begin{minipage}{.85\textwidth}\sffamily\lite\small\RaggedRight}%
% 	{\end{minipage}\end{Sbox}%
% 		\begin{center}\setlength{\fboxsep}{6pt}%
% 		\colorbox{litegrey}{\TheSbox}\end{center}}

\newenvironment{codesnippet}{%
	\begin{Sbox}\begin{minipage}{\textwidth}\sffamily\small}%
	{\end{minipage}\end{Sbox}%
		\begin{center}%
		\colorbox{litegrey}{\TheSbox}\end{center}}



\usepackage{fancyhdr}
\pagestyle{fancy}

%\renewcommand{\chaptermark}[1]{\markboth{#1}{}}
\renewcommand{\sectionmark}[1]{\markright{\thesection\ - #1}}

\fancyhf{}

\fancyhead[LO]{Sección \rightmark} % \thesection\ 
\fancyfoot[LO]{\small{Nombre Apellido, Nombre Apellido, Nombre Apellido}}
\fancyfoot[RO]{\thepage}
\renewcommand{\headrulewidth}{0.5pt}
\renewcommand{\footrulewidth}{0.5pt}
\setlength{\hoffset}{-0.8in}
\setlength{\textwidth}{16cm}
%\setlength{\hoffset}{-1.1cm}
%\setlength{\textwidth}{16cm}
\setlength{\headsep}{0.5cm}
\setlength{\textheight}{25cm}
\setlength{\voffset}{-0.7in}
\setlength{\headwidth}{\textwidth}
\setlength{\headheight}{13.1pt}

\renewcommand{\baselinestretch}{1.1}  % line spacing


% \setcounter{secnumdepth}{2}
\usepackage{underscore}
\usepackage{caratula}
\usepackage{url}


% ******************************************************** %
%              TEMPLATE DE INFORME ORGA2 v0.1              %
% ******************************************************** %
% ******************************************************** %
%                                                          %
% ALGUNOS PAQUETES REQUERIDOS (EN UBUNTU):                 %
% ========================================
%                                                          %
% texlive-latex-base                                       %
% texlive-latex-recommended                                %
% texlive-fonts-recommended                                %
% texlive-latex-extra?                                     %
% texlive-lang-spanish (en ubuntu 13.10)                   %
% ******************************************************** %



\begin{document}


\thispagestyle{empty}
\materia{Organización del Computador II}
\submateria{Primer Cuatrimestre de 2014}
\titulo{Trabajo Práctico II}
\subtitulo{subtitulo del trabajo}
\integrante{Tomas Shaurli}{671/10}{tshaurli@gmail.com}
\integrante{Sebastian Aversa}{379/11}{sebastianaversa@gmail.com}
\integrante{Fernando Gabriel Otero}{424/11}{fergabot@gmail.com}

\maketitle
\newpage

\thispagestyle{empty}
\vfill
\begin{abstract}
En el presente trabajo se describe el an\'alisis de siertos filtros para im\'agenes, comparando las versiones realizadas en C y en assembler con procesamiento SIMD
\end{abstract}
\newpage
\thispagestyle{empty}
\vspace{3cm}
\tableofcontents
\newpage


\normalsize
\newpage

\section{Objetivos generales}

El objetivo de este Trabajo Pr\'actico fue el aprendizaje de las instrucciones SIMD(Single Instruction, Multiple Data) mediante su uso en filtros de im\'agenes, y las comparaciones entre los algoritmos en assembler realizados con \'estas instrucciones y los mismos algoritmos en C.
Tanto en los lenguajes assembler como C se compil\'o y ensambl\'o para arquitecturas Intel x64.

\section{Contexto}

Los algoritmos descriptos y analizados fueron probados en las computadoras de los laboratorios de la facultad. Todas las comparaciones fueron realizadas en las mismas computadoras para un correcto an\'alisis.
Las implementaciones en assembler fueron hechas específicamente para el uso de las instrucciones SIMD de Intel x64 con los siguientes par\'ametros:
\begin{itemize}
	\item -felf64:
	\item -g:
	\item -F:
	\item dwarf:
\end{itemize}
Los sigientes par\'ametros fueron usados en el compilador de C:
\begin{itemize}
	\item ggdb:
	\item -Wall:
	\item -std=c99:
	\item -pedantic:
	\item -m64:
	
	\item pkg-config:
	\item --cflags:
	\item --libs opencv:
\end{itemize}


%\newpage
%\input{introduccion}
\newpage
\section{Filtro de tiles}
\subsection{Enunciado}

 Programar el filtro \textit{tiles} en lenguaje C y luego en ASM haciendo uso de las instrucciones vectoriales (\textbf{SSE}).

\vspace*{0.3cm} \noindent
\textbf{Experimento 1 - análisis el código generado}

Utilizar la herramienta \verb|objdump| para verificar como el compilador de C deja ensamblado el código C. Como es el código generado, ¿cómo se manipulan las variables locales?¿le parece que ese código generado podría optimizarse?

\vspace*{0.3cm} \noindent
\textbf{Experimento 2 - optimizaciones del compilador}

Compile el código de C con optimizaciones del compilador, por ejemplo, pasando el flag \verb|-O1|\footnote{agregando este flag a \texttt{CCFLAGS64} en el makefile}. 
¿Qué optimizaciones realizó el compilador?
¿Qué otros flags de optimización brinda el compilador?
¿Para qué sirven?


\vspace*{0.3cm} \noindent
\textbf{Experimento 3 - secuencial vs. vectorial}

	Realice una medición de las diferencias de performance entre las versiones
	de C y ASM (el primero con -O1, -O2 y -O3).\\
	¿Como realizó la medición?¿Cómo sabe que su medición es una buena medida?¿Cómo afecta a la medición la existencia de \emph{outliers}\footnote{en español, valor atípico: \url{http://es.wikipedia.org/wiki/Valor_atípico}}?¿De qué manera puede minimizar su impacto?¿Qué resultados obtiene si mientras corre los tests ejecuta otras aplicaciones que utilicen al máximo la CPU? 
	Realizar un análisis \textbf{riguroso} de los resultados y acompañar con un gráfico que presente estas diferencias.


\vspace*{0.3cm} \noindent
\textbf{Experimento 4 - cpu vs. bus de memoria}

	Se desea conocer cual es el mayor limitante a la performance de este filtro en su versión ASM.

	¿Cuál es el factor que limita la performance en este caso? En caso de que el limitante
	fuera la intensidad de cómputo, entonces podrían agregarse instrucciones que realicen
	accesos a memoria y la performance casi no debería sufrir. La inversa puede aplicarse
	si el limitante es la cantidad de accesos a memoria.
	
	Realizar un experimento, agregando múltiples instrucciones de un mismo tipo y realizar un análisis
	del resultado. Acompañar con un gráfico.

\vspace*{0.3cm} \noindent
\textbf{Experimento 5 (\textit{opcional}) - secuencial vs. vectorial (parte II)}

	Si vemos a los pixeles como una tira muy larga de bytes, este filtro en
	realidad no requiere ningún procesamiento de datos en paralelo. Esto podría
	significar que la velocidad del filtro de C puede aumentarse hasta casi
	alcanzar la del de ASM. ¿ocurre esto?
	
	Modificar el filtro para que en vez de acceder a los bytes de a uno a la vez
	se accedan como tiras de 64 bits y analizar la performance.
	
\newpage

\subsection{An\'alisis previo}
El filtro de tiles define un rect\'angulo en una imagen en color. La funci\'on debe replicar dicho sector de la im\'agen en el resto de la imagen original.
Para esto, se nos ocurrieron tres formas de encararlo: con la f\'ormula dada en el enunciado (versi\'on C), guardando el primer elemento del rect\'angulo a copiar y guardando un puntero a ese primer elemento.

Formula:

i y j representan fila y columna respectivamente.

dst(i,j) = src [(i mod tamy)+offsety][(j mod tamx)+offsetx]

\subsection{Implementaci\'on en C}

La implementaci\'on en C del ejercicio tiles es tomar un rect\'angulo de la imagen original y repartirla por la imagen destino. Se resuelve con dos ciclos for anidados usando la forma dada en el enuciado.

\subsection{Implementaci\'on en assembler}
Para la resoluci\'on en assembler creamos los dos ciclos anidados con los registros de prop\'osito general. AL comienzo del segundo siclo se hace la copia.
\begin{codesnippet}
\begin{verbatim}
	MOVDQU XMM0, [RDI]
	MOVDQU [RSI], XMM0
\end{verbatim}
\end{codesnippet}
Luego, para ciclar y recorrer toda la matriz se modifican los valores de RSI y RDI cuando es necesario.
\subsection{Resultado de los experimentos}

\vspace*{0.3cm} \noindent
\textbf{Experimento 1 - análisis el código generado}
El compilador genera más variables locales que las usadas en nuestro c\'odigo assembler. En ese sentido, estamos realizando una mejora al c\'odigo de memoria y velocidad


\vspace*{0.3cm} \noindent
\textbf{Experimento 2 - optimizaciones del compilador}
El compilador ofrece los flags -O1 a -03 para distintos niveles de optimizaci\'on, asi como los flags para indicar si es assembler de intel, amd, etc. Esto ultimo hace m\'as espec\'ifico el codigo assembler y correr\'ia mejor en una maquina con el microprocesador seleccionado,aunque no correr\'ia en uno de otra marca o de una familia de microprocesadores anterior a la elegida.

\newpage
\section{Filtro Popart}
\subsection{Enunciado}

\subsection*{Filtro \textit{Popart}}

  Programar el filtro \textit{Popart} en lenguaje C y en en ASM haciendo uso de 
  las instrucciones vectoriales (\textbf{SSE}).

\vspace*{0.3cm} \noindent
\textbf{Experimento 1 - saltos condicionales}

	Se desea conocer que tanto impactan los saltos condicionales
	en el código del ejercicio anterior con \verb|-O1|.\\
	Para poder medir esto, una posibilidad es quitar las comparaciones
	al procesar cada pixel. Por más que la imagen resultante no sea correcta,
	será posible tomar una medida del impacto de los saltos condicionales.
	Analizar como varía la performance. 
	
	Si se le ocurren, mencionar otras posibles formas de medir el impacto de los saltos condicionales.

\vspace*{0.3cm} \noindent
\textbf{Experimento 2 - cpu vs. bus de memoria}

	¿Cuál es el factor que limita la performance en este caso? 
	
	Realizar un experimento, agregando múltiples instrucciones de un mismo tipo
	y realizar un análisis 	del resultado. Acompañar con un gráfico.


\vspace*{0.3cm} \noindent
\textbf{Experimento 3 - prefetch}

  La técnica de \textit{prefetch} es otra forma de optimización que puede
  realizarse. Su sustento teórico es el siguiente:
  
  Suponga un algoritmo que en cada iteración tarda n ciclos en obtener un dato y una cantidad
  similar en procesarlo. Si el algoritmo lee el dato $i$ y luego lo procesa,
  desperdiciará siempre n ciclos esperando entre que el dato llega y que se comienza
  a procesar efectivamente. Un algoritmo más inteligente podría pedir el 
  dato $i+1$ al comienzo del ciclo de proceso del dato $i$ (siempre suponiendo
  que el dato $i$ pidió en la iteración $i-1$. De esta manera, a la vez que el
  procesador computa todas las instrucciones de la iteración $i$, se estáran trayendo
  los datos de la siguiente iteración, y cuando esta última comience, los datos ya
  habrán llegado.

  \vspace*{0.2cm}
  Estudiar esta técnica y proponer una aplicación al código del filtro en la versión ASM.
  Programarla y analizar el resultado. ¿Vale la pena hacer prefetching?

\vspace*{0.3cm} \noindent
\textbf{Experimento 3 - secuencial vs. vectorial}

  Analizar cuales son las diferencias de performace entre las versiones de C y ASM. 
  Realizar gráficos que representen estas diferencias.



\subsection{Implementaci\'on en C}
Para la implementaci\'on en C usamos un vector con los posibles reemplazos para cada pixel, y una funci\'on que devuelve cual va a ser ese reemplazo mediante comparaciones.
\begin{codesnippet}
\begin{verbatim}

rgb_t colores[] = { {255,   0,   0},
                    {127,   0, 127},
                    {255,   0, 255},
                    {  0,   0, 255},
                    {  0, 255, 255} };

			
int pop_art(unsigned int r, unsigned int g, unsigned int b)
{
	unsigned int suma = r + g + b;
	int res;
	
	if 		 (suma < 153)	{res = 0;}
	else if (suma < 306)	{res = 1;}
	else if (suma < 459)	{res = 2;}
	else if (suma < 612)	{res = 3;}
	else 					{res = 4;}
	
	return res;
}
\end{verbatim}
\end{codesnippet}

\subsection{Implementaci\'on en assembler}
Armamos máscaras para las comparaci\'ones con suma(i,j).
\begin{codesnippet}
\begin{verbatim}
seisDoce:          DD 0x264, 0x264, 0x264, 0x264    ; 612, 612, 612, 612			
cuatroCincoNueve:  DD 0x1CB, 0x1CB, 0x1CB, 0x1CB    ; 459, 459, 459, 459	
tresOSeis:         DD 0x132, 0x132, 0x132, 0x132    ; 306, 306, 306, 306	
unoCincoTres:      DD 0x99, 0x99, 0x99, 0x99        ; 153, 153, 153, 153
\end{verbatim}
\end{codesnippet}

Tambi\'en para reemplazar los pixeles (solo cito el primer caso).
\begin{codesnippet}
\begin{verbatim}
colores0:          DB  0xFF, 0x0,  0x0, 0x0,        ; 255,   0,   0,   0
                   DB  0xFF, 0x0,  0x0, 0x0,        ; 255,   0,   0,   0
                   DB  0xFF, 0x0,  0x0, 0x0,        ; 255,   0,   0,   0
                   DB  0xFF, 0x0,  0x0, 0x0         ; 255,   0,   0,   0

\end{verbatim}
\end{codesnippet}
El resto de las m\'ascaras, usadas con \verb|PSHUFB|.\\ son:

\begin{itemize}
	\item pixel0BGR : Para tener un pixel (R,G y B) en cada \textit{double word}. El cuarto Byte se rellena con ceros
			
	\item soloPrimero : Lo usamos para trabajar sobre el primer Byte de cada \textit{double word}
		
	\item pixelFinal : Sirve para sacar el ceros ubicado en el cuarto Byte de cada \textit{double word}. El ultimo DB se rellena con ceros
\end{itemize}
\'Estas m\'ascaras son usadas en \textit{pop_art} para discriminar los Bytes de R, G y B y guardar la suma de cada uno en un \textit{double word}

\subsection{Resultado de los experimentos}
%\begin{figure}
%  \begin{center}
%	\includegraphics[scale=0.66]{imagenes/logouba.jpg}
%	\caption{Logo}
%	\label{Logo}
%  \end{center}
%\end{figure}



\newpage
\section{Filtro Temperature}
\subsection{Enunciado}

\subsection*{Filtro \textit{Temperature}}

  Programar el filtro \textit{Temperature} en lenguaje C y en en ASM haciendo uso de 
  las instrucciones vectoriales (\textbf{SSE}).


\vspace*{0.3cm} \noindent
\textbf{Experimento 1}

  Analizar cuales son las diferencias de performace entre las versiones de C y ASM. 
  Realizar gráficos que representen estas diferencias.
  
\subsection{An\'alisis previo}

\subsection{Implementaci\'on en C}
\paragraph{\textbf{Titulo del parrafo} } Bla bla bla bla.
Esto se muestra en la figura~\ref{nombreparareferenciar}.
\subsection{Implementaci\'on en assembler}
\begin{codesnippet}
\begin{verbatim}

struct Pepe {

    ...

};
\end{verbatim}
\end{codesnippet}
\subsection{Resultado de los experimentos}
%\begin{figure}
%  \begin{center}
%	\includegraphics[scale=0.66]{imagenes/logouba.jpg}
%	\caption{Logo}
%	\label{Logo}
%  \end{center}
%\end{figure}



\newpage
\section{Filtro LDR}
\subsection{Enunciado}

\subsection*{Filtro \textit{LDR}}
  Programar el filtro \textit{LDR} en lenguaje C y en
  ASM haciendo uso de las instrucciones \textbf{SSE}.

\vspace*{0.3cm} \noindent
\textbf{Experimento 1}

  Analizar cuales son las diferencias de performace entre las versiones de C y ASM. 
  Realizar gráficos que representen estas diferencias.
  
\subsection{An\'alisis previo}

\subsection{Implementaci\'on en C}
\paragraph{\textbf{Titulo del parrafo} } Bla bla bla bla.
Esto se muestra en la figura~\ref{nombreparareferenciar}.
\subsection{Implementaci\'on en assembler}
\begin{codesnippet}
\begin{verbatim}

struct Pepe {

    ...

};
\end{verbatim}
\end{codesnippet}
\subsection{Resultado de los experimentos}
%\begin{figure}
%  \begin{center}
%	\includegraphics[scale=0.66]{imagenes/logouba.jpg}
%	\caption{Logo}
%	\label{Logo}
%  \end{center}
%\end{figure}



\newpage
\section{Conclusiones y trabajo futuro}
asdf4


\end{document}

